\documentclass{homework}
% \usepackage{graphicx} % Required for inserting images
\usepackage{amsmath}
\usepackage{gensymb}
\usepackage{subfig}
\usepackage{float}

\class{APPM 5560}
\title{Homework 4}
\author{Sean Jennings}
\date{February 26, 2024}

\begin{document}
\maketitle

\section{Problem 1}
\begin{letterenum}
\item Each state is recurrent, since the whole state space is finite and closed
(trivially). It is also irredicible since each state communicates with
each other state. Specifically, we have
\begin{enumerate}
    \item $3\to1\to2\to3$ and 
    \item $3\to4\to5\to3$.
\end{enumerate}
So each state $x$ communicates with state 3 and state 3 communicates with each
state $y$, therefore $x\to y$.

\item Each state in the set $C = \{5, 6, 7\}$ is recurrent, since $C$ is
finite, closed ($p(x, y) = 0$ for all $x\in C$, $y \notin C$) and irreducible 
($p(x, y) > 0$ for all $x,y \in C$).

The set of states $T = \{1, 2, 3, 4\}$ are all transient. For each 
$z \in\{1, 2\}$ and $y\in C$ we have $z\to3\to4\to y$, so each $x \in T$
we have $\rho_{xy} > 0$ since $x$ communicates with $y$. But since $C$
is closed, $\rho_{yx} = 0<1$. Therefore $\rho_{xx} < 1$ for each $x \in T$:
there is a nonzero probability that the state will not return to $T$.
\end{letterenum}

\section{Problem 2}
Since the stationary distribution $\pi \in \R^{1 \times 3}$ satisfies
\begin{equation*}
    \pi = \pi p,
\end{equation*}
we have
\begin{equation*}
    \pi (p - I) = 0.
\end{equation*}
Taking the transpose yields
\begin{equation*}
    (p^T - I) \pi^T = 0.
\end{equation*}
So to find the stationary distribution is in the null space of the matrix
\begin{equation*}
    p^T - I = \mat{
        -2/3 & 0 & 1 \\
        1/3 & -1/2 & 0 \\
        1/3 & 1/2 & -1
    }
\end{equation*}
Converting to a system of equations and solving algebraically yields the 
relations:
\begin{equation*}
    \pi_3 = 2/3 \pi_1 \qquad \pi_2 = 2/3 \pi_1
\end{equation*}
So $\pi$ is of the form
\begin{equation*}
    \pi = \mat{1 & 2/3 & 2/3} \pi_1.
\end{equation*}
The elements of $\pi$ must sum to 1, so we have $\pi_1 = 3/7$ and
\begin{equation*}
    \pi = \mat{3/7 & 2/7 & 2/7}
\end{equation*}

\section{Problem 3}
\begin{letterenum}
\item Suppose the set $D$ is closed, finite and have all its states are reccurent.
This does {\bf not} imply it is irreducible. As a counter example, consider
the whole state space $S$ of two-state markov chain with probability transition matrix
\begin{equation*}
    p = \mat{1 & 0 \\ 0 & 1}.
\end{equation*}
$S$ is trivially closed and clearly finite. Futhermore, since each state
$x \in S$ always transitions to itself, we have $\rho_{xx} = 1$
and $x$ is reccurent.  However, the states do not communicate with each other
and therefore $S$ is reducible.

\item A irreducible, finite set $D$ with all recurrent states is not necessarily
closed. For example, take $D = \{1\}$ in the two-state chain with prob. transition
matrix
\begin{equation*}
    p = \mat{1/2 & 1/2 \\ 1/2 & 1/2}
\end{equation*}
State 1 is recurrent since $1 \to 2 \to 1$. The set $D$ is also irreducible
since each state in $D$ communicates with each other state 
in $D$ (this is just $1\to 1$). However, $D$ is clearly not closed since
$p(1, 2) = 1/2 > 0$ and $1\in D$, $2\notin D$.

\item A irreducible, closed set $D$ with all states recurrent is not necessarily
finite. A state space of a symmetric 2D random walk is an example of such an 
infinite set. Each state in the random walk communicates with each other and 
the state space is trivially closed. Futhermore, each state is recurrent
since the expected number of returns to a given state $x$ is an inifinite
series which diverges.

\end{letterenum}


\section{Problem 4}
\begin{letterenum}
\item Each state in $R = \{5, 6, 7, 8\}$ is recurrent, while each
state of $T = \{1, 2, 3, 4\}$ is transient.

\item For $\rho_{11}$ the only path for state 1 to return to itself
is from 1 to 2 to 3 to 4, then back to 1. So
\begin{equation*}
    \rho_{11} = p(1,2)p(2,3)p(3,4)p(4,1) = a^4
\end{equation*}
Similarly, there is one path from 1 to 3, so
\begin{equation*}
    \rho_{13} = p(1, 2)p(2, 3) = a^2
\end{equation*}
Since state 5 always transitions to itself, we have
\begin{equation*}
    \rho_{51} = 0 \qquad \rho_{55} = 1
\end{equation*}

\item The expected times state 1 will be visited from state 2 is
\begin{equation*}
    E_2[N_1] = \frac{\rho_{21}}{1 - \rho_{11}} = \frac{a^3}{1 - a^4}
\end{equation*}

\item The desired probability is
\begin{align*}
    P_1(T_2^k < \infty) &= P(\text{reaches to from 1})
        \cdot P(\text{returns to 2 from 2})^{k-1} \\
        &= \rho_{12} \rho_{22}^{k-1} \\
        &= a (a^4)^{k-1} \\
        &= a^{4k-3}
\end{align*}




\end{letterenum}

\section{Problem 5}

\begin{letterenum}
\item Take some initial distribution $q_0\in \R^{1\times 2}$ such that
$q_0(1) + q_0(2) = 1$. Then, the distribution at the next timestep
is computed by right multiplying by the probability transition matrix $p$:
\begin{equation}
    q_{n+1} = q_n p = q_n \mat{
        1 - a & a \\
        b & 1 - b
    }
    \label{eq:dynamics}
\end{equation}
As a system of scalar equations, this is equivalent to:
\begin{align*}
    q_{n+1}(1) &= (1 - a) q_n(1) + a q_n(2) \\
    q_{n+1}(2) &= b q_n(1) + (1-b) q_n(2)
\end{align*}

\item Let $\alpha_n \in \R^{1 \times 2}$ be the row vector representing
the difference from the stationary distribution:
\begin{equation}
    \alpha_n = q_n - \pi
    \label{eq:change-of-var}
\end{equation}

Combining equations (\ref{eq:dynamics}) and (\ref{eq:change-of-var}),
we have
\begin{align*}
    \alpha_{n+1} &= q_{n+1} - \pi \\
        &= q_n p - \pi \\
        &= (\alpha_n + \pi) p - \pi \\
        &= \alpha_n p + \pi (p - I).
\end{align*}
where $I$ is the 2x2 identity matrix.
We see that the term contain the stationary distribution is 0, since
\begin{align*}
    \pi (p - I) &= \frac{1}{a + b} \mat{b & a} \mat{
        -a & a \\
        b & -b
    } \\
    &= \frac{1}{a + b} \mat{ab - ba & ab - ba} = 0
\end{align*}
so that the iteration in the change of variables is actually equivalent
to the original iteration:
\begin{equation}
    \alpha_{n+1} = \alpha_n p
    \label{eq:dynamics-cov}
\end{equation}
Taking $x_n, y_n \in \R$ to be the elements of $\alpha_n = \mat{x_n & y_n}$
we can also rewrite equation (\ref{eq:dynamics-cov}) as the scalar equations:
\begin{align*}
    x_{n+1} &= (1 - a) x_n + a y_n \\
    y_{n+1} &= b x_n + (1-b) y_n
\end{align*}

\item To show that $x_n$ and $y_n$ exponentially decay to zero, we 
will show that the 1-norm of $\alpha_n$ has an exponential bound of the
form:
\begin{equation*}
    |x_n| + |y_n| = \normone{\alpha_n} \leq C e^{-\lambda n}
\end{equation*}
for some $C > 0$ and $\lambda > 0$. Clearly, this implies that both
$x_n$ and $y_n$ decay to zero, since the preceding inequality implies
(along with the fact: $|z| \geq 0$ for all $z\in\R$) that
\begin{equation*}
    |x_n| \leq C e^{-\lambda n} - |y_n| \leq C e^{-\lambda n}
\end{equation*}
and likewise for $|y_n|$.


We begin by writing equation (\ref{eq:dynamics-cov}) in terms of $\alpha_0$.
Since each succeeding value is a left multiplication by $p$, we end up
with the $n$-th power of $p$:
\begin{equation*}
    \alpha_n = \alpha_0 p^n.
\end{equation*}
We take the transpose since working with column vectors is more familiar:
\begin{equation}
    \alpha_n^T = (p^T)^n \alpha_0^T
    \label{eq:dynamics-power}
\end{equation}

We note that $p^T$ is diagonalizable and can be rewritten as $p^T = ADA^{-1}$
with $A,D \in\R^{2\times 2}$ ($D$ diagonal) given by
\begin{equation*}
    A = \mat{
        1 & b \\
        -1 & a
    } \qquad
    D = \mat{
        1 - a - b & 0 \\
        0 & 1
    } \qquad
    A^{-1} = \frac{1}{a + b} \mat{
        a & -b \\
        1 & 1
    }
\end{equation*}
By the associative property of matrix multiplication,
interior terms of the power cancel
\begin{equation*}
    (p^T)^n = (ADA^{-1})(ADA^{-1})\cdots(ADA^{-1}) = AD^nA^{-1}.
\end{equation*}
Substituting into equation (\ref{eq:dynamics-power}) yields:
\begin{equation}
    \alpha_n^T = A \mat{
        (1 - a - b)^n & 0 \\
        0 & 1
    } A^{-1} \alpha_0^T.
    \label{eq:dyanmics-eigenvec-basis}
\end{equation}
Since the elements of $q_n$ and $\pi$ sum to one, we have
\begin{equation*}
    x_n + y_n = \sum_i \alpha_n(i) = \sum_i (q_n(i) - \pi(i))
        = \sum_i q_n(i) - \sum_i \pi(i)
        = 1 - 1 = 0.
\end{equation*}

This implies that the component (corresponding to the eigenvalue of 1)
of $A^{-1} \alpha^T$
(which is the represention of $\alpha^T$ in the basis of the eigenvectors) is zero,
for
\begin{equation*}
    A^{-1} \alpha_0^T = \frac{1}{a + b} \mat{
        ax_0 - by_0 \\
        x_0 + y_0
    } = \frac{1}{a + b} \mat{
        ax_0 - by_0 \\
        0
    }.
\end{equation*}

Thus, equation (\ref{eq:dyanmics-eigenvec-basis}) reduces to
\begin{align*}
    \alpha_n^T &= A \mat{
        (1 - a - b)^n & 0 \\
        0 & 1
    } \frac{1}{a + b} \mat{ax_0 - by_0 \\ 0} \\
    &= \frac{(1 - a - b)^n}{a + b} \mat{
        1 & b \\
        -1 & a
    }\mat{
        ax_0 - by_0 \\ 0
    } \\
    &= \frac{(1 - a - b)^n}{a + b} \mat{
        ax_0 - by_0 \\
        by_0 - ax_0
    } \\
    &= \frac{(1 - a - b)^n}{a + b} \mat{
        a & -b \\
        -a & b
    } \alpha_0^T
\end{align*}

When $a + b \neq 1$, taking the 1-norm of both sides yields
the desired exponential bound:
\begin{align*}
    \normone{\alpha_n^T} &= \mnormone{
        \frac{(1 - a - b)^n}{a + b} \mat{
            a & -b \\
            -a & b
        } \alpha_0^T } \\
    &\leq \biggl|\frac{(1 - a - b)^n}{a + b}\biggr|
        \mnormone{\mat{
            a & -b \\
            -a & b
        }} \normone{\alpha_0} \\
    &= \frac{|1 - a - b|^n}{a + b} \max\{2a, 2b\} \normone{\alpha_0} \\
    &= C e^{-\lambda n}
\end{align*}

where $C$ is given by:
\begin{equation*}
    C = \frac{2(|x_0| + |y_0|)\max\{a, b\}}{a + b}
\end{equation*}
and $\lambda$ by:
\begin{equation*}
    \lambda = \ln \biggl(\frac{1}{|1 - a - b|}\biggr).
\end{equation*}
We have $C>0$ for any $\alpha_0 \neq 0$, so it remains to show that $\lambda > 0$.
Since $0 < a < 1$ and $0 < b < 1$, we have $-1 < (1 - a - b) < 1$.
so that $1/|1 - a - b| > 1$. Thus, $\lambda > 0$ and exponential decay
towards the stationary distribution is proven.

Note that $\lambda$ is well-defined only under the assumption that $a + b \neq 1$.
If $a + b = 1$, then we see from the previous equations that $\alpha_n$
is exactly zero for all $n>0$. That is, the markov chain converges exactly to
the stationary distribution after the first step.










\end{letterenum}

\section{Problem 6}
\begin{letterenum}
\item The probability of leaving each odd state must sum to 1, so we have
\begin{equation*}
    b + c = 1.
\end{equation*}
The probability of leaving each even state must also sum to 1, so
\begin{equation*}
    a + c + d = 1
\end{equation*}

\item Since the chain is symmetric between even and odd states, the
function which returns the state modulo 2 is a homomorphism. Consider
the two-state markov chain with state $Q_n$ which is 0 if the state of
the original markov chain $X_n$ is even and 1 is $X_n$ is odd.

Then the probability transition matrix for this reduced markov chain is given 
by
\begin{equation*}
    p = \mat{
        d & a + c \\
        b + c & 0
    } = \mat{
        d & 1 - d \\
        1 & 0
    }
\end{equation*}
Put $a = 1 - d$ and $b = 1$ into the results for the stationary distribution
from the previous exercise to see that $\pi_Q$, the stationary distribution for the
reduced markov chain, is
\begin{equation*}
    \pi_Q = \mat{\pi_{even} & \pi_{odd}} = \mat{\frac{1}{2-d} & \frac{1}{2-d}}.
\end{equation*}

Since the original chain is finite and irreducible (given the constraints on a,b,c, and d),
the stationary distribution must be unique. We reason that the ratio between
even and odd states must be equal to that of the reduced chain, multiplied by 
a normalization factor $\lambda$. Let $\pi \in \R^{1 \times 10}$ be the
stationary distribution for the original chain. Then,
\begin{equation*}
    \pi(i) = \biggl\{
        \begin{array}{ll}
            \lambda / (2 - d) &\quad \text{if }i \text{ is even} \\
            \lambda (1 - d) / (2 - d) &\quad \text{if }i \text{ is odd}
        \end{array}
    \biggr.
\end{equation*}
Since the elements of $\pi$ sum to 1, we have
\begin{align*}
    1 = \sum_{i=1}^{10} \pi(i) &= \sum_{i\in \text{Even}} \pi(i)
        + \sum_{i\in \text{Odd}} \pi(i) \\
        &= 5 \biggl(\frac{\lambda}{2 - d} + \frac{\lambda (1 - d)}{2-d} \biggr) \\
        &= 5\lambda.
\end{align*}
So $\lambda = 1/5$ and 
\begin{equation*}
    \pi(i) = \biggl\{
        \begin{array}{ll}
            \frac{1}{5(2 - d)} &\quad \text{if }i \text{ is even} \\
            \frac{1 - d}{5(2 - d)} &\quad \text{if }i \text{ is odd}
        \end{array}
    \biggr.
\end{equation*}
\end{letterenum}

\section{Problem 7}
\begin{letterenum}
\item Let $T$ be the number of steps until the markov chain reaches
either state 4 or 5. Starting from state 1, the chain can reach
states 4 or 5 via both state 2 and 3, so we start by computed the expectation
of $T$ starting from these intermediate states.

Starting from state 3, $T$ is a geometric R.V. with prob. of success
$p(3,4) + p(3, 5) = 3/4$. Therefore,
\begin{equation*}
    E_3[T] = 1 / (p(3,4 + p(3,5))) = 4/3
\end{equation*}

Starting from state 2, we use the total expectation theorem to compute
the expectation based on the probability of each possible step:
\begin{align*}
    E_2[T] &= (1 + E_2[T]) p(2, 2) + (1 + E_3[T]) p(2, 3) + 1 p(2,4) \\
        &= E_2[T] p(2, 2) + E_3[T] p(2, 3) + (p(2, 2) + p(2, 3) + p(2, 4)) \\
    (1 - p(2, 2)) E_2[T] &= E_3[T] p(2, 3) + 1 \\
    E_2[T] &= \frac{E_3[T] p(2, 3) + 1}{1 - p(2, 2)} \\
        &= \frac{4/3 \cdot 1/8 + 1}{1 - 1/2} = 2 (1 + 1/6) \\
        &= 7/3
\end{align*}

Starting from state 1, the number of steps to reach state 2 is geometric
with success 1/10. So the expected number of steps to reach state 2 is 10.
Since all paths must go through state 2, the expected number of steps 
to reach 4 or 5 from 1 is 10 more than the expectation starting from 2:
\begin{equation*}
    E_1[T] = 10 + E_2[T] = 37/3
\end{equation*}

\item To compute the probability, $\rho_{14}$, that the chain 
starting from state 1 reaches state 4, we similarly work backwards from state 3.

By the total probability law, and because $\rho_{54} = 0$, we have
\begin{align*}
    \rho_{34} &= p(3, 4) + p(3, 3) \rho_{34} \\
    (1 - p(3, 3)) \rho_{34} &= p(3, 4) \\
    \rho_{34} &= \frac{p(3,4)}{1 - p(3, 3)} = \frac{1/4}{1 - 1/4} \\
        &= 1/3 
\end{align*}

Similarly, from state 2, we have
\begin{align*}
    \rho_{24} &= \sum_{i=1}^{5} P_2(T_4 < \infty | X_1 = i) P_2(X_1 = i) \\
        &= \sum_{i=1}^5 \rho_{i4} p(2, i) \\
        &= \rho_{24} p(2, 2) + \rho_{34} p(2, 3) + \rho_{44} p(2, 4) \\
    (1 - p(2, 2))\rho_{24} &= \rho_{34} p(2, 3) + \rho_{44} p(2, 4) \\
    \rho_{24} &= \frac{\rho_{34} p(2, 3) + \rho_{44} p(2, 4)}{1 - p(2, 2)} \\
        &= \frac{1/3 \cdot 1/8 + 1/8}{1/2} \\
        &= 2 (4 / 24) \\
        &= 1/3
\end{align*}

The probability of reaching 4 from state 1 is the same as from state 2
since:
\begin{align*}
    \rho_{14} &= p(1, 1) \rho_{14} + p(1, 2) \rho_{12} \\
    (1 - p(1, 1)) \rho_{14} &= p(1, 2) \rho_{12}
\end{align*}
and $p(1, 1) + p(1, 2) = 1$. Thus,
\begin{equation*}
    \rho_{14} = \rho_{24} = 1/3
\end{equation*}

\end{letterenum}

\section{Problem 8}
Let the state $X_n$ be the distance between the spider and the fly
at each step $n$.
We define the state space of the markov chain to be $S = \Z_{\geq 0} \cup \{\infty\}$
where $\Z_{\geq 0}$ is the set of all non-negative integers and the state
$\infty$ indicates that fly has escaped.

Since the spider always moves towards the fly, the distance (excluding when $\infty$
the fly escapes) is non-increasing. So for a given initial state $x_0\in\Z_{\geq 0}$,
the set $C(x_0) = \{0, 1, \dots, x_0\} \cup \{\infty\}$ is closed.

The transition probabilities for $\infty$ and 0 are:
\begin{equation*}
    p(0, 0) = p(\infty, \infty) = 1
\end{equation*}
Otherwise, when $x\in \Z_{\geq 0}$ and $x > 1$, the transition probabilities
are given by
\begin{equation*}
    p(x, y) = \biggl\{
        \begin{array}{ll}
            0.3 &\quad \text{if }y = x \\
            0.4 &\quad \text{if }y = x - 1 \\
            0.3 &\quad \text{if }y = x - 2
        \end{array}
    \biggr.
\end{equation*}
When $x=1$, the probabilities are conditioned on the possibility of escape:
\begin{equation*}
    p(1, y) = \biggl\{
        \begin{array}{ll}
            0.5 &\quad \text{if }y = \infty \\
            0.5 \cdot (0.3 + 0.3) &\quad \text{if }y = 1 \\
            0.5 \cdot 0.4 &\quad \text{if }y = 0
        \end{array}
    \biggr.
\end{equation*}
Given that the fly does not escape, there is a 0.6 probability that the fly will
move so that the distance remains 1. Thus, the total probability that $y=1$
is:

$P$(fly moves 1 square) $= P($fly moves 1 square$|$fly does not escape) $P($fly does not escape$)$.

\begin{letterenum}
\item The probability that the fly escapes starting from state 2 is $\rho_{2\infty}$.
The total probability law gives the following equations:
\begin{align*}
    \rho_{2\infty} &= p(2, 2) \rho_{2\infty} + p(2, 1) \rho_{1\infty}
        = 0.3 \rho_{2\infty} + 0.4 \rho_{1\infty} \\
    \rho_{1\infty} &= p(1, 1) \rho_{1\infty} + p(1, \infty) \rho_{\infty\infty}
        = 0.3 \rho_{1\infty} + 0.5
\end{align*}
which admits the solution:
\begin{align*}
    \rho_{1\infty} &= 5/7 \\
    \rho_{2\infty} &= 4/7 \rho_{1\infty} = \boxed{20/49 \approx 0.408}
\end{align*}

\item Let $T = \min\{n: X_n \in\{0, \infty\}\}$. Similarly to the previous
problem, we use the total expectation theorem. From state 1, we have
\begin{align*}
    E_1[T] &= p(1, 1) (1 + E_1[T]) + p(1, 0) 1 + p(1, \infty) 1 \\
        &= (p(1, 1) + p(1, 0) + p(1, \infty)) + p(1, 1) E_1[T] \\
        &= 1 + p(1, 1) E_1[T]  = \frac{1}{1 - p(1, 1)} \\
        &= \frac{1}{1 - 0.3} = 10/7
\end{align*}

So from state 2:
\begin{align*}
    E_2[T] &= p(2, 0) 1 + p(2, 1) ( 1 + E_1[T]) + p(2, 2) (1 + E_2[T]) \\
        &= 1 + p(2, 1) E_1[T] + p(2, 2) E_2[T] \\
        &= \frac{1 + p(2, 1) E_1[T]}{1 - p(2, 2)} \\
        &= \frac{1 + 0.4 \cdot 10/7}{0.7} \\
        &= \frac{11/7}{0.7} \\
        &= \boxed{110/49 \approx 2.245}
\end{align*}

\item The expected number of times in imminent danger is given by
\begin{align*}
    E_2[N] &= \sum_{k=1}^\infty P_2(N_1 \geq k) \\
        &= \sum_{k=1}^\infty \rho_{21} \rho_{11}^{k-1} \\
        &= \frac{\rho_{21}}{1 - \rho_{11}}
\end{align*}

We have
\begin{equation*}
    \rho_{11} = p(1, 1) = 0.3
\end{equation*}
and
\begin{equation*}
    \rho_{21} = p(2, 2) \rho_{21} + p(2, 1) = \frac{p(2, 1)}{1 - p(2, 2)}
        = \frac{0.4}{1 - 0.3} = 4/7
\end{equation*}
So the expected number of times staying in danger is:
\begin{equation*}
    E_2[N] = \frac{4/7}{1 - 0.3} = \frac{4/7}{7/10} = \boxed{40/49 \approx 0.816}
\end{equation*}

    
\end{letterenum}








\end{document}
